% Generated draft: requires mathematical and source review.
\documentclass[11pt]{article}
\usepackage[margin=25mm]{geometry}
\usepackage{fontspec}
\setmainfont{Latin Modern Roman}
\usepackage{amsmath,amssymb,mathtools,unicode-math}
\setmathfont{Latin Modern Math}
\usepackage{enumitem,needspace,xurl,hyperref}
\setlength{\parindent}{0pt}
\setlength{\parskip}{5pt}
\setlength{\emergencystretch}{3em}
\newcommand{\N}{\mathbb{N}}
\newcommand{\Z}{\mathbb{Z}}
\newcommand{\Q}{\mathbb{Q}}
\newcommand{\R}{\mathbb{R}}
\newcommand{\C}{\mathbb{C}}
\newcommand{\F}{\mathbb{F}}
\newcommand{\A}{\mathbb{A}}
\newcommand{\PP}{\mathbb{P}}
\newcommand{\E}{\mathbb{E}}
\newcommand{\eps}{\varepsilon}
\newcommand{\dd}{\,\mathrm{d}}
\begin{document}
\section*{UnsolvedMath ID: 30000717}
\textbf{Source problem code:} catalogue entry
\subsection{Definitions and mathematical statement}
Let \(A=\R[X_1,\ldots,X_n]\), let \(f\in A\), and write \(f_i=\partial f/\partial X_i\). Put \(f^*=\inf_{x\in\R^n}f(x)\), without an assumption that the infimum is attained. The cleaned statement defines
\[
J=J_{\mathrm{skew}}=(X_jf_i-X_if_j:1\leq i<j\leq n).
\]
Let \(V_\R(J)\) be its common real zero set, let \(\Sigma A^2\) be the finite sums of polynomial squares, and let \(\sqrt J=\{h\in A:h^r\in J\text{ for some }r\geq1\}\) be the ordinary algebraic radical, not silently the real radical.

The cleaned source asks whether the following four conditions are equivalent for every indicated polynomial, with no attained-minimum hypothesis:
\[
\text{(i)}\quad f(x)\geq0\quad(x\in\R^n),
\]
\[
\text{(ii)}\quad f(x)\geq0\quad(x\in V_\R(J)),
\]
\[
\text{(iii)}\quad f\in\Sigma A^2+\sqrt J,
\]
\[
\text{(iv)}\quad\forall\varepsilon>0,\quad f+\varepsilon\in\Sigma A^2+J.
\]
In (iv), the sum-of-squares representation and ideal multipliers may depend on \(\varepsilon\); no common degree bound is requested. The exact zero-level certificate in (iii) and the positive-perturbation certificates in (iv) are separate clauses. The source does not explicitly require \(f^*\) finite, so no bounded-below hypothesis is inserted beyond the proposed conditions themselves.

There is a material generator discrepancy: the original wording describes the ideal generated by \(X_jf_i\) and \(X_if_j\) separately, whereas the cleaned version uses their differences. The literal separate-product ideal would be
\[
J_{\mathrm{sep}}=(X_jf_i,\ X_if_j:1\leq i<j\leq n),
\]
which need not equal \(J_{\mathrm{skew}}\). For example, for \(f=X_1^2+X_2^2\), the first difference vanishes identically, while the separate products generate \((X_1X_2)\). This diagnoses inequivalent readings, not a counterexample to the proposed four-way statement. The supplied source-verification claim does not explain the discrepancy, so the exact original ideal must be reconciled rather than silently chosen.

The supplied literature triage checked 2026-08-21 attributes the replacement conjecture to Schweighofer and reports generic and special-case progress but no unrestricted resolution in its assessment. No exact partial-result hypotheses or proof of any proposed equivalence is supplied.
\subsection{Short English statement}
The curated question proposes four equivalent nonnegativity and sum-of-squares conditions after replacing the gradient ideal by the displayed tangency ideal, without assuming a minimum is attained. The original text instead names separate-product generators, a substantive discrepancy requiring source reconciliation.
\subsection{Sources}
\begin{itemize}
\item Local UnsolvedMath registry, permanent ID 30000717, catalogue entry, \emph{Gradient-Tentacle Certificates for Polynomial Nonnegativity}; the supplied assessment attributes the conjecture to Schweighofer.
\item \emph{Reelle Algebraische Geometrie} (2007), Oberwolfach Reports, pp. 753--820, DOI \url{https://doi.org/10.4171/owr/2007/14}.
\item Supplied reference ``gradient tentacles paper'': \url{https://arxiv.org/abs/math/0511617}.
\item Supplied verification claims comparison with the official report; supplied literature triage checked 2026-08-21 does not reconcile the two generator descriptions or state exact special-case hypotheses.
\end{itemize}
\end{document}
